\scanpage{55}
\mspage{42}
% archive sheet number: 54

\deleted{$I$ finie} l'ens $W_{I^{\unsure{\circ|0}}}$, et $C_I$ (il vaut mieux sans doute le noter aussi $C_{I^\circ}$) \unsure{défini|définis} au \S{} précédent. Il s'agit de vérifier les axiomes. L'axiome $\mathrm{Cof}_0$ est immédiat, et bien connu pour $I$ une catégorie quelconque.

$\mathrm{Cof}_1$ : $C_{I^\circ}$ contient les iso, c'est clair, parce que $C$ les contient. Stable par composition $X_{\unsure{\cdot|\bullet}}\xrightarrow{u}Y_{\cdot}\xrightarrow{v}Z_{\cdot}$, \unsure{regardons} $K\subset J$, posons $w=\unsure{v\circ u|vu}$

\diagram{\begin{tikzcd}[ampersand replacement=\&,column sep=small,row sep=small]
 X_K \arrow[r] \arrow[d,dash] \& X_{\deleted{\unsure{K}}J} \arrow[d] \& \& \\
 Y_K \arrow[r] \arrow[d] \& Y_{K,J} \arrow[r,"{u_{K,J}}"] \arrow[d] \& Y_J \arrow[d,dash] \& \\
 Z_K \arrow[r] \& \cdot \arrow[r,"{u'_{K,J}}"] \arrow[rr,phantom,shift right=2ex,"{\underbrace{\hspace{8em}}_{w_{K,J}}}"'] \& Z_{K,J} \arrow[r,"{v_{K,J}}"] \& Z_J
\end{tikzcd}}

$w_{K,J}$ est composé de $v_{K,J}$, qui est dans $C$, et $u'_{K,J}$, qui l'est car se déduit de $u_{K,J}$ par \unsure{cochgt} de base. (NB \unsure{On} a déjà fait le raisonnement avec un seul $u$, et \uline{trois} ouverts $J\subset K\subset L$. Il
